\documentclass{article}
\usepackage{amsmath, amssymb, amsthm}
\usepackage{hyperref}
\usepackage{geometry}
\geometry{margin=1in}

\title{Erd\H{o}s Problem \#409}
\date{Last updated: 2025-08-31}
\author{Source: erdosproblems.com}

\begin{document}
\maketitle

\noindent\textbf{Status:} OPEN \quad \textbf{No prize}

\noindent\textbf{Tags:} number theory, iterated functions

\medskip
\noindent\textbf{Problem Statement:}

How many iterations of $n\mapsto \phi(n)+1$ are needed before a prime is reached? Can infinitely many $n$ reach the same prime? What is the density of $n$ which reach any fixed prime?

\bigskip
\noindent\textbf{Remarks:}

A problem of Finucane. One can also ask similar questions about $n\mapsto \sigma(n)-1$: do iterates of this always reach a prime? If so, how soon? (It is easily seen that iterates of this cannot reach the same prime infinitely often, since they are non-decreasing.)

This problem is somewhat ambiguously phrased. Let $F(n)$ count the number of iterations of $n\mapsto \phi(n)+1$ before reaching a prime. The number of iterations required is A039651 in the OEIS.

Cambie notes in the comments that $F(n)=o(n)$ is trivial, and $F(n)=1$ infinitely often. Presumably the intended question is to find 'good' upper bounds for $F(n)$.

This is discussed in problem B41 of Guy's collection \cite{Gu04}.

\bigskip
\noindent\textbf{References:}

[Gu04] Guy, Richard K., Unsolved problems in number theory. (2004), xviii+437.

\bigskip
\noindent\small{Source: \url{https://www.erdosproblems.com/409}}
\end{document}
